\section*{不定积分的概念与性质}

1. 选择题

设 \(F(x)\) 为 \(f(x)\) 在某区间 \(I\) 内的一个原函数，则下列命题不成立的是（C）。

(A) \(\frac{d}{dx}\left(\int f(x)\,dx\right)=f(x)\)（正确，这是微积分基本定理）

(B) \(\int F'(x)\,dx=F(x)+C\)（正确，这是原函数的定义）

(C) \(\int F'(x)\,dx=f(x)+C\)（错误，应为 \(F(x)+C\)）

(D) \(d\left(\int f(x)\,dx\right)=f(x)\,dx\)（正确，这是微分的定义）

2. 判断题

(1) 若 \(f(x)\) 的某个原函数为常数，则 \(f(x)=0\)。（正确）

(2) 一切初等函数在其有定义的区间上都有原函数。（正确）

(3) 一切初等函数在其定义域内都有原函数。（错误，可能存在孤立点导致问题）

(4) 函数 \(f(x)\) 的无穷多个原函数组成的集合称为 \(f(x)\) 的不定积分。（正确）

3. 求下列不定积分。

(1) \(\int\left(3x^{3}-5x^{2}+\frac{3}{x^{2}}\right)dx\)。

解：
\[ \int\left(3x^{3}-5x^{2}+\frac{3}{x^{2}}\right)dx=\frac{3}{4}x^{4}-\frac{5}{3}x^{3}-\frac{3}{x}+C. \]

(2) \(\int\sqrt{x\sqrt{x\sqrt{x}}}\,dx\)。

解：先化简被积函数：
\[ \sqrt{x\sqrt{x\sqrt{x}}}=\sqrt{x\cdot x^{1/2}\cdot x^{1/4}}=\sqrt{x^{1+1/2+1/4}}=\sqrt{x^{7/4}}=x^{7/8}. \]
因此
\[ \int x^{7/8}\,dx=\frac{8}{15}x^{15/8}+C. \]

(3) \(\int(2\tan x+3\cot x)^{2}\,dx\)。

解：先展开并化简：
\[ (2\tan x+3\cot x)^{2}=4\tan^{2}x+12\tan x\cot x+9\cot^{2}x=4\tan^{2}x+12+9\cot^{2}x, \]
\[ =4(\sec^{2}x-1)+12+9(\csc^{2}x-1)=4\sec^{2}x+9\csc^{2}x-1. \]
因此
\[ \int(2\tan x+3\cot x)^{2}\,dx=\int(4\sec^{2}x+9\csc^{2}x-1)\,dx=4\tan x-9\cot x-x+C. \]

(4) \(\int e^{3x}\left(3^{x}-e^{-2x}\right)dx\)。

解：
\[ \int e^{3x}\left(3^{x}-e^{-2x}\right)dx=\int e^{3x}3^{x}\,dx-\int e^{3x}e^{-2x}\,dx \]
\[ =\int e^{x(3+\ln 3)}\,dx-\int e^{x}\,dx=\frac{1}{3+\ln 3}e^{x(3+\ln 3)}-e^{x}+C. \]

(5) \(\int\left(\frac{1}{x}-\frac{3}{\sqrt{1-x^{2}}}\right)dx\)。

解：
\[ \int\left(\frac{1}{x}-\frac{3}{\sqrt{1-x^{2}}}\right)dx=\ln|x|-3\arcsin x+C. \]

(6) \(\int\frac{\sqrt{x}-2\sqrt[3]{x^{2}}+1}{4\sqrt{x}}\,dx\)。

解：先化简被积函数：
\[ \frac{\sqrt{x}-2\sqrt[3]{x^{2}}+1}{4\sqrt{x}}=\frac{1}{4}-\frac{1}{2}x^{1/6}+\frac{1}{4}x^{-1/2}. \]
因此
\[ \int\frac{\sqrt{x}-2\sqrt[3]{x^{2}}+1}{4\sqrt{x}}\,dx=\frac{1}{4}x-\frac{3}{7}x^{7/6}+\frac{1}{2}x^{1/2}+C. \]

(7) \(\int\frac{2^{x-1}-5^{x-1}}{10^{x}}\,dx\)。

解：首先简化被积函数：
\[ \frac{2^{x-1}-5^{x-1}}{10^{x}}=\frac{\frac{2^{x}}{2}-\frac{5^{x}}{5}}{2^{x}\cdot 5^{x}}=\frac{1}{2}\cdot\frac{1}{5^{x}}-\frac{1}{5}\cdot\frac{1}{2^{x}}=\frac{1}{2}\cdot 5^{-x}-\frac{1}{5}\cdot 2^{-x}. \]
所以
\[ \int\frac{2^{x-1}-5^{x-1}}{10^{x}}\,dx=\int\left(\frac{1}{2}\cdot 5^{-x}-\frac{1}{5}\cdot 2^{-x}\right)dx. \]
分别积分：
\[ \int\frac{1}{2}\cdot 5^{-x}\,dx=\frac{1}{2}\int 5^{-x}\,dx=\frac{1}{2}\left(-\frac{5^{-x}}{\ln 5}\right)=-\frac{5^{-x}}{2\ln 5}, \]
\[ \int\frac{1}{5}\cdot 2^{-x}\,dx=\frac{1}{5}\int 2^{-x}\,dx=\frac{1}{5}\left(-\frac{2^{-x}}{\ln 2}\right)=-\frac{2^{-x}}{5\ln 2}. \]
因此
\[ \int\left(\frac{1}{2}\cdot 5^{-x}-\frac{1}{5}\cdot 2^{-x}\right)dx=-\frac{5^{-x}}{2\ln 5}-\left(-\frac{2^{-x}}{5\ln 2}\right)=\frac{2^{-x}}{5\ln 2}-\frac{5^{-x}}{2\ln 5}. \]
加上积分常数，最终结果为
\[ \int\frac{2^{x-1}-5^{x-1}}{10^{x}}\,dx=\frac{2^{-x}}{5\ln 2}-\frac{5^{-x}}{2\ln 5}+C. \]

(8) \(\int\frac{(1-x)^{2}}{x(1+x^{2})}\,dx\)。

解：部分分式分解：
\[ \frac{(1-x)^{2}}{x(1+x^{2})}=\frac{1}{x}+\frac{-2}{1+x^{2}}. \]
故
\[ \int\frac{(1-x)^{2}}{x(1+x^{2})}\,dx=\ln|x|-2\arctan x+C. \]

(9) \(\int\frac{x^{2}}{1+x^{2}}\,dx\)。

解：
\[ \int\frac{x^{2}}{1+x^{2}}\,dx=\int\left(1-\frac{1}{1+x^{2}}\right)dx=x-\arctan x+C. \]

(10) \(\int\frac{1+\cos^{2}2x}{1-\cos 2x}\,dx\)。

解：先用 \(\cos 2x=1-2\sin^{2}x\) 化简。
\[ \frac{1+\cos^{2}2x}{1-\cos 2x}=\frac{2-4\sin^{2}x+4\sin^{4}x}{2\sin^{2}x}=\csc^{2}x-2+2\sin^{2}x. \]
因此
\[ \int\frac{1+\cos^{2}2x}{1-\cos 2x}\,dx=\int(\csc^{2}x-2+2\sin^{2}x)\,dx \]
\[ =-\cot x-2x+\int(1-\cos 2x)\,dx=-\cot x-2x+x-\frac{1}{2}\sin 2x+C \]
\[ =-\cot x-x-\frac{1}{2}\sin 2x+C. \]

4. 曲线方程问题

已知曲线上任一点 \((x,y)\) 的斜率为 \(\frac{2}{\sqrt{1-x^{2}}}\)，且曲线过点 \(\left(\frac{1}{2},0\right)\)。试求曲线方程。

证明：由题意得微分方程
\[ \frac{dy}{dx}=\frac{2}{\sqrt{1-x^{2}}}. \]
两边积分得
\[ y=2\arcsin x+C. \]
代入点 \(\left(\frac{1}{2},0\right)\)，\(\arcsin\frac{1}{2}=\frac{\pi}{6}\)，于是
\[ 0=2\cdot\frac{\pi}{6}+C=\frac{\pi}{3}+C\ \Rightarrow\ C=-\frac{\pi}{3}. \]
因此曲线方程为
\[ y=2\arcsin x-\frac{\pi}{3}. \]

5. 质点位移问题

设某质点的速率（单位：米/秒）为 \(v(t)=3t^{2}-1\)，且当 \(t=1\) 秒时位移为 \(10\) 米，求 \(t=3\) 秒时的位移。

证明：位移函数为速度的一个原函数：
\[ s(t)=\int v(t)\,dt=\int(3t^{2}-1)\,dt=t^{3}-t+C. \]
由 \(s(1)=10\) 得
\[ 1-1+C=10\ \Rightarrow\ C=10. \]
故 \(s(t)=t^{3}-t+10\)，因此
\[ s(3)=27-3+10=34, \]
即 \(t=3\) 秒时的位移为 34 米。

6. 证明题

证明：若 \(\int f(x)\,dx=F(x)+C\)，则 \(\int f(ax+b)\,dx=\frac{1}{a}F(ax+b)+C\)。

证明：令 \(u=ax+b\)，则 \(du=a\,dx\)，故 \(dx=\frac{1}{a}du\)。于是
\[ \int f(ax+b)\,dx=\int f(u)\cdot\frac{1}{a}\,du=\frac{1}{a}\int f(u)\,du=\frac{1}{a}F(u)+C=\frac{1}{a}F(ax+b)+C. \]
这完成了证明。
